\chapter{Materiais e Métodos}
\label{cap:metodologia}

<Descreva o circuito e o procedimento de simulação. Todo o fluxo roda em
Docker via \texttt{./claw-spice}: geração do circuito, simulação LTspice,
extração de medidas e renderização do esquemático.>

\section{Circuito}
\label{sec:circuito}

<Descreva a topologia e os valores dos componentes. O esquemático da
Figura~\ref{fig:esquematico} foi renderizado a partir do arquivo
\texttt{.asc} com \texttt{./claw-spice show <circuito>.asc --png}.>

\begin{figure}[htb]
  \centering
  % Substitua pelo PNG do esquemático em imagens/generated/
  %\includegraphics[width=0.8\textwidth]{<circuito>.png}
  \fbox{\parbox[c][4cm][c]{0.8\textwidth}{\centering <esquemático:
  imagens/generated/<circuito>.png>}}
  \caption{Esquemático do circuito simulado.}
  \label{fig:esquematico}
\end{figure}

\begin{table}[htb]
  \centering
  \caption{Componentes do circuito.}
  \label{tab:componentes}
  \begin{tabular}{lll}
    \toprule
    Referência & Componente & Valor \\
    \midrule
    R1 & Resistor  & \SI{1}{\kilo\ohm} \\
    C1 & Capacitor & \SI{1}{\micro\farad} \\
    V1 & Fonte de tensão & degrau de \SI{0}{\volt} a \SI{5}{\volt} \\
    \bottomrule
  \end{tabular}
\end{table}

\section{Simulação}
\label{sec:simulacao}

<Descreva a análise usada (\texttt{.tran}, \texttt{.ac}, \texttt{.op}),
o passo/janela de tempo e as diretivas \texttt{.meas} definidas. Liste os
comandos executados, por exemplo:>

\begin{verbatim}
./claw-spice sim run <circuito>.cir
./claw-spice log summary <circuito>.log
./claw-spice raw plot <circuito>.raw V(out) \
    --output latex/imagens/generated/<plot>.svg --png
\end{verbatim}

<A netlist completa está no Apêndice~\ref{ap:netlist}.>
